% ============================================================================
% Soluciones del Módulo 2
% ============================================================================

\begin{solucion}{m2-1}
Filtramos con dos condiciones, elegimos columnas y ordenamos; después
resumimos el resultado con \codigo{summarise()}.

\begin{rcode}
library(tidyverse)
\end{rcode}

\begin{rcode}
ventas <- read_csv("datasets/ventas_retail.csv", show_col_types = FALSE)

tienda3_desc <- ventas %>%
  filter(tienda_id == 3, descuento_pct >= 10) %>%   # ambas condiciones
  select(fecha, producto_id, descuento_pct, total) %>%
  arrange(desc(total))
tienda3_desc

tienda3_desc %>%
  summarise(ventas = n(), total = sum(total))
\end{rcode}
\begin{rsalida}
# A tibble: 4 x 4
  fecha      producto_id descuento_pct total
  <date>           <dbl>         <dbl> <dbl>
1 2023-11-05          27            10 1439.
2 2023-03-02          26            20 1227.
3 2023-12-20          17            10 1098.
4 2023-07-08           6            10  742.
# A tibble: 1 x 2
  ventas total
   <int> <dbl>
1      4 4506.
\end{rsalida}

La tienda 3 hizo 4 ventas con descuento de 10\% o más, que suman
\$4,506; la mayor fue de \$1,439 el 5 de noviembre. Observa que en
\codigo{summarise(total = sum(total))} el nombre nuevo puede coincidir con
el de la columna: dentro de la función, \codigo{total} todavía es la
columna original.
\end{solucion}

\begin{solucion}{m2-2}
Combinamos tres condiciones en \codigo{filter()} (el producto activo y el
stock bajo el mínimo) y calculamos el faltante con \codigo{mutate()}.

\begin{rcode}
library(tidyverse)
productos <- read_csv("datasets/productos.csv", show_col_types = FALSE)

resurtir <- productos %>%
  filter(activo, stock_actual < stock_minimo) %>%
  mutate(faltante = stock_minimo - stock_actual) %>%
  select(nombre_producto, proveedor, stock_actual, stock_minimo,
         faltante) %>%
  arrange(desc(faltante))
resurtir
resurtir %>% count(proveedor, wt = faltante, sort = TRUE)
\end{rcode}
\begin{rsalida}
# A tibble: 3 x 5
  nombre_producto proveedor   stock_actual stock_minimo faltante
  <chr>           <chr>              <dbl>        <dbl>    <dbl>
1 Producto 03     Proveedor 3           22           26        4
2 Producto 33     Proveedor 1           42           45        3
3 Producto 06     Proveedor 1           36           38        2
# A tibble: 2 x 2
  proveedor       n
  <chr>       <dbl>
1 Proveedor 1     5
2 Proveedor 3     4
\end{rsalida}

Solo 3 productos activos están bajo su mínimo. Por unidades faltantes hay
que llamar primero al Proveedor 1 (5 unidades en dos productos), aunque el
producto con el mayor faltante individual (Producto 03, 4 unidades) es del
Proveedor 3. \codigo{filter(activo, \ldots)} funciona porque
\codigo{activo} ya es lógica.
\end{solucion}

\begin{solucion}{m2-3}
\codigo{case\_when()} evalúa las reglas en orden, así que basta con dar el
límite superior de cada rango. El porcentaje se calcula después de
\codigo{count()}.

\begin{rcode}
library(tidyverse)
clientes <- read_csv("datasets/clientes.csv", show_col_types = FALSE)

clientes %>%
  mutate(rango_edad = case_when(
    edad < 30 ~ "18-29",
    edad < 45 ~ "30-44",
    edad < 60 ~ "45-59",
    .default  = "60+"
  )) %>%
  count(rango_edad) %>%
  mutate(pct = round(n / sum(n) * 100, 1))
\end{rcode}
\begin{rsalida}
# A tibble: 4 x 3
  rango_edad     n   pct
  <chr>      <int> <dbl>
1 18-29         33  16.5
2 30-44         52  26  
3 45-59         51  25.5
4 60+           64  32  
\end{rsalida}

El grupo más grande es el de 60 años o más (64 clientes, 32\%) y el más
chico el de 18 a 29 (33 clientes, 16.5\%): una base de clientes madura.
\end{solucion}

\begin{solucion}{m2-4}
El promedio de una columna lógica es la proporción de \codigo{TRUE}: con
\codigo{mean(es\_premium)} obtenemos directamente el porcentaje de premium.

\begin{rcode}
library(tidyverse)
clientes <- read_csv("datasets/clientes.csv", show_col_types = FALSE)

perfil_ciudad <- clientes %>%
  group_by(ciudad) %>%
  summarise(clientes    = n(),
            pct_premium = round(mean(es_premium) * 100, 1),
            edad_prom   = round(mean(edad), 1),
            .groups = "drop") %>%
  arrange(desc(clientes))
perfil_ciudad
perfil_ciudad %>% slice_max(pct_premium, n = 1)
\end{rcode}
\begin{rsalida}
# A tibble: 10 x 4
   ciudad      clientes pct_premium edad_prom
   <chr>          <int>       <dbl>     <dbl>
 1 Puebla            28        42.9      51.3
 2 CDMX              27        48.1      56.2
 3 Guadalajara       22        27.3      49.2
 4 Cancún            20        30        45  
 5 Toluca            20        30        51  
 6 León              18        16.7      39.1
 7 Mérida            17        23.5      47.6
 8 Querétaro         17        35.3      53.7
 9 Tijuana           17        41.2      37.8
10 Monterrey         14        71.4      51.1
# A tibble: 1 x 4
  ciudad    clientes pct_premium edad_prom
  <chr>        <int>       <dbl>     <dbl>
1 Monterrey       14        71.4      51.1
\end{rsalida}

Puebla y la CDMX tienen más clientes, pero Monterrey tiene la mayor
proporción de premium: 71.4\% de sus 14 clientes. Es una plaza chica pero de
alto valor.
\end{solucion}

\begin{solucion}{m2-5}
\begin{rcode}
library(tidyverse)
transacciones <- read_csv("datasets/transacciones.csv",
                          show_col_types = FALSE)

transacciones %>%
  mutate(dia = wday(fecha, label = TRUE, abbr = FALSE)) %>%
  group_by(dia) %>%
  summarise(transacciones = n(),
            ingresos = sum(total_transaccion),
            ticket_promedio = round(mean(total_transaccion), 2),
            .groups = "drop") %>%
  mutate(participacion_pct = round(ingresos / sum(ingresos) * 100, 1))
\end{rcode}
\begin{rsalida}
# A tibble: 7 x 5
  dia       transacciones ingresos ticket_promedio participacion_pct
  <ord>             <int>    <dbl>           <dbl>             <dbl>
1 domingo             146  148703.           1019.              15.5
2 lunes               136  128878.            948.              13.4
3 martes              129  119364.            925.              12.4
4 miércoles           155  141329.            912.              14.7
5 jueves              146  145280.            995.              15.1
6 viernes             158  161266.           1021.              16.8
7 sábado              130  115172.            886.              12  
\end{rsalida}

El viernes es el día con más transacciones (158) e ingresos (\$161,266,
16.8\% del total), seguido del domingo; el sábado es el más flojo (12\%).
Conviene reforzar el personal los viernes y domingos.
\end{solucion}

\begin{solucion}{m2-6}
\begin{rcode}
library(tidyverse)
empleados <- read_csv("datasets/empleados.csv", show_col_types = FALSE)
fecha_corte <- as.Date("2024-01-01")

empleados_antig <- empleados %>%
  mutate(
    antig_calc = time_length(interval(fecha_ingreso, fecha_corte),
                             "years"),
    diferencia = abs(antig_calc - antiguedad_anos),
    grupo = case_when(
      antig_calc < 3  ~ "1. Menos de 3 años",
      antig_calc <= 6 ~ "2. 3 a 6 años",
      .default        = "3. Más de 6 años"
    )
  )
max(empleados_antig$diferencia)          # diferencia máxima

empleados_antig %>%
  group_by(grupo) %>%
  summarise(empleados = n(),
            salario_prom = round(mean(salario_mensual), 2),
            .groups = "drop")
\end{rcode}
\begin{rsalida}
[1] 0.05245902
# A tibble: 3 x 3
  grupo              empleados salario_prom
  <chr>                  <int>        <dbl>
1 1. Menos de 3 años        41       27354.
2 2. 3 a 6 años             34       27821.
3 3. Más de 6 años          25       27228.
\end{rsalida}

La diferencia máxima entre ambos cálculos es de 0.05 años (unos 19
días): la columna original dividía los días entre 365, mientras que
\codigo{time\_length()} considera los años bisiestos. El salario promedio
casi no cambia con la antigüedad (entre \$27,228 y \$27,821): en estos datos
simulados la antigüedad no se refleja en el sueldo, algo que en una empresa
real merecería una revisión de política salarial.
\end{solucion}

\begin{solucion}{m2-7}
Solo necesitamos del catálogo la categoría y el costo, así que los
seleccionamos dentro del \emph{join}.

\begin{rcode}
library(tidyverse)
transacciones <- read_csv("datasets/transacciones.csv",
                          show_col_types = FALSE)
productos <- read_csv("datasets/productos.csv", show_col_types = FALSE)

margen_cat <- transacciones %>%
  left_join(select(productos, producto_id, categoria, costo),
            by = "producto_id") %>%
  mutate(utilidad = (precio_venta - costo) * cantidad) %>%
  group_by(categoria) %>%
  summarise(con_perdida = sum(utilidad < 0),   # antes de resumir
            ingresos = sum(total_transaccion),
            utilidad = sum(utilidad),
            .groups = "drop") %>%
  mutate(margen_pct = round(utilidad / ingresos * 100, 1)) %>%
  arrange(desc(utilidad))
margen_cat
\end{rcode}
\begin{rsalida}
# A tibble: 10 x 5
   categoria   con_perdida ingresos utilidad margen_pct
   <chr>             <int>    <dbl>    <dbl>      <dbl>
 1 Alimentos            13  121543.   73870.       60.8
 2 Electrónica          27  132672.   73067.       55.1
 3 Juguetes             30  125093.   57714.       46.1
 4 Belleza              66  182143.   56054.       30.8
 5 Automotriz           26  104085.   55538.       53.4
 6 Deportes             20   86171.   48301.       56.1
 7 Mascotas             10   57668.   35423.       61.4
 8 Libros               21   57162.   25757.       45.1
 9 Hogar                24   63407.   13079.       20.6
10 Ropa                 15   30048.    5835.       19.4
\end{rsalida}

Alimentos es la categoría que más utilidad deja (\$73,870), aunque
Belleza tiene los mayores ingresos (\$182,143). Belleza es también la más
problemática: 66 transacciones con utilidad negativa y un margen de 30.8\%.
Ropa y Hogar tienen los márgenes más bajos (19--21\%). Cuidado con el orden
dentro de \codigo{summarise()}: si calculas \codigo{utilidad = sum(utilidad)}
\emph{antes} de \codigo{sum(utilidad < 0)}, la segunda usaría la suma ya
calculada y el conteo saldría 0.
\end{solucion}

\begin{solucion}{m2-8}
Trabajamos con listas de clientes \emph{únicos} de cada tabla
(\codigo{distinct()}) para que los conteos sean de clientes y no de ventas.

\begin{rcode}
library(tidyverse)
ventas <- read_csv("datasets/ventas_retail.csv", show_col_types = FALSE)
transacciones <- read_csv("datasets/transacciones.csv",
                          show_col_types = FALSE)
clientes <- read_csv("datasets/clientes.csv", show_col_types = FALSE)

cli_ventas <- ventas %>% distinct(cliente_id)
cli_trans  <- transacciones %>% distinct(cliente_id)

# (a) En ambas tablas
cli_ventas %>% semi_join(cli_trans, by = "cliente_id") %>% nrow()
# (b) En transacciones pero no en ventas
cli_trans %>% anti_join(cli_ventas, by = "cliente_id") %>% nrow()
# (c) Registrados que no aparecen en ninguna
compradores <- bind_rows(cli_ventas, cli_trans) %>% distinct()
clientes %>% anti_join(compradores, by = "cliente_id") %>%
  select(cliente_id, nombre, fecha_registro)
\end{rcode}
\begin{rsalida}
[1] 170
[1] 28
# A tibble: 0 x 3
# i 3 variables: cliente_id <dbl>, nombre <chr>,
#   fecha_registro <date>
\end{rsalida}

170 clientes compraron en ambas tablas y 28 aparecen solo en
transacciones. Ningún cliente registrado quedó sin comprar en alguna de las
dos: el resultado de (c) es una tabla vacía (0 filas), lo cual también es
una respuesta válida.
\end{solucion}

\begin{solucion}{m2-9}
Calculamos el porcentaje en formato largo (agrupado por tienda) y al final
pivotamos a ancho.

\begin{rcode}
library(tidyverse)
transacciones <- read_csv("datasets/transacciones.csv",
                          show_col_types = FALSE)
tiendas <- read_csv("datasets/tiendas.csv", show_col_types = FALSE)

pagos_tienda <- transacciones %>%
  left_join(select(tiendas, tienda_id, nombre_tienda), by = "tienda_id") %>%
  group_by(nombre_tienda, metodo_pago) %>%
  summarise(ingresos = sum(total_transaccion), .groups = "drop_last") %>%
  mutate(pct = round(ingresos / sum(ingresos) * 100, 1)) %>%
  ungroup() %>%
  select(-ingresos) %>%
  pivot_wider(names_from = metodo_pago, values_from = pct,
              values_fill = 0) %>%
  arrange(desc(Efectivo))
pagos_tienda
\end{rcode}
\begin{rsalida}
# A tibble: 10 x 5
   nombre_tienda    Efectivo `Tarjeta Crédito` `Tarjeta Débito`
   <chr>               <dbl>             <dbl>            <dbl>
 1 Sucursal Oeste       45.5              35.6             10.7
 2 Sucursal Este        37.5              27               25.6
 3 Sucursal Norte       37.1              43.5             14.2
 4 Sucursal Plaza A     36.2              21.5             31.4
 5 Sucursal Plaza B     33.5              42.8             18.1
 6 Sucursal Sur         29.3              36.4             29.6
 7 Sucursal Outlet      29.1              29.5             28  
 8 Sucursal Centro      28.9              44               21.2
 9 Sucursal Premium     24.6              40.8             25.2
10 Sucursal Express     18.8              47.6             20.8
# i 1 more variable: Transferencia <dbl>
\end{rsalida}

La Sucursal Oeste es la más dependiente del efectivo: 45.5\% de sus
ingresos. En el otro extremo, la Sucursal Express cobra solo 18.8\% en
efectivo y 47.6\% con tarjeta de crédito. Como los porcentajes se calcularon
con el resultado agrupado por tienda (\codigo{.groups = "drop\_last"}),
cada fila suma 100.
\end{solucion}

\begin{solucion}{m2-10}
Primero completamos los meses faltantes; si no, \codigo{lag()} compararía
contra el mes equivocado. \codigo{group\_by(tienda\_id)} hace que
\codigo{lag()} y \codigo{cumsum()} reinicien en cada tienda.

\begin{rcode}
library(tidyverse)
ventas <- read_csv("datasets/ventas_retail.csv", show_col_types = FALSE)

tienda_mes <- ventas %>%
  group_by(tienda_id, mes) %>%
  summarise(ingresos = sum(total), .groups = "drop") %>%
  complete(tienda_id, mes, fill = list(ingresos = 0)) %>%
  arrange(tienda_id, mes) %>%
  group_by(tienda_id) %>%
  mutate(cambio_pesos = ingresos - lag(ingresos),
         acumulado    = cumsum(ingresos)) %>%
  ungroup()
nrow(tienda_mes)

tienda_mes %>%
  slice_min(cambio_pesos, n = 3) %>%        # las 3 mayores caídas
  mutate(across(where(is.numeric), ~ round(.x, 2)))
\end{rcode}
\begin{rsalida}
[1] 120
# A tibble: 3 x 5
  tienda_id mes     ingresos cambio_pesos acumulado
      <dbl> <chr>      <dbl>        <dbl>     <dbl>
1         8 2023-08     781.      -12373.    50452 
2         8 2023-03    1222.       -9457.    17226.
3         2 2023-05    3698.       -7626.    25458.
\end{rsalida}

La tabla completa tiene 120 filas (10 tiendas $\times$ 12 meses) gracias
a \codigo{complete()}. La mayor caída fue la de la tienda 8 en agosto:
\$12,373 menos que en julio, al vender solo \$781 en el mes.
\end{solucion}

\begin{solucion}{m2-11}
Seguimos el mismo orden que en la sección~\ref{sec:m2-limpieza}:
duplicados, texto, números y fechas, y al final reglas de negocio.

\begin{rcode}
library(tidyverse)
lista_proveedor <- tibble(
  id        = c(1, 2, 2, 3, 4, 5, 6),
  categoria = c("electronica", "Electrónica ", "Electrónica ",
                "HOGAR", "hogar", "Ropa", " ropa"),
  precio    = c("$1,250.00", "$899.50", "$899.50", "450", "$0",
                "$320.75", "-15"),
  fecha     = c("2023-05-01", "02/05/2023", "02/05/2023", "2023-05-03",
                "04/05/2023", "2023-05-05", "06/05/2023")
)

lista_limpia <- lista_proveedor %>%
  distinct() %>%                                   # 1. duplicados
  mutate(
    categoria = str_to_title(str_squish(categoria)),   # 2. texto
    categoria = recode(categoria, "Electronica" = "Electrónica"),
    precio    = parse_number(precio),                  # 3. número
    fecha     = as_date(parse_date_time(fecha,         # 4. fecha
                                        orders = c("ymd", "dmy"))),
    sku       = str_c("SKU-", str_pad(id, 4, pad = "0"))
  ) %>%
  filter(precio > 0)                               # 5. regla de negocio
lista_limpia

lista_limpia %>%
  group_by(categoria) %>%
  summarise(productos = n(), precio_prom = mean(precio),
            .groups = "drop")
\end{rcode}
\begin{rsalida}
# A tibble: 4 x 5
     id categoria   precio fecha      sku     
  <dbl> <chr>        <dbl> <date>     <chr>   
1     1 Electrónica  1250  2023-05-01 SKU-0001
2     2 Electrónica   900. 2023-05-02 SKU-0002
3     3 Hogar         450  2023-05-03 SKU-0003
4     5 Ropa          321. 2023-05-05 SKU-0005
# A tibble: 3 x 3
  categoria   productos precio_prom
  <chr>           <int>       <dbl>
1 Electrónica         2       1075.
2 Hogar               1        450 
3 Ropa                1        321.
\end{rsalida}

De 7 registros quedaron 4: se eliminaron el duplicado del id 2, el precio
en cero (id 4) y el negativo (id 6). Las categorías quedaron unificadas y
las fechas en un solo formato.
\end{solucion}

\begin{solucion}{m2-12}
Construimos una tabla maestra, la validamos con \codigo{stopifnot()},
calculamos las tres hojas y las exportamos con \paquete{writexl}.

\begin{rcode}
library(tidyverse)
library(writexl)
transacciones <- read_csv("datasets/transacciones.csv",
                          show_col_types = FALSE)
productos <- read_csv("datasets/productos.csv", show_col_types = FALSE)
tiendas   <- read_csv("datasets/tiendas.csv", show_col_types = FALSE)

maestra_t <- transacciones %>%
  left_join(select(productos, producto_id, nombre_producto, categoria),
            by = "producto_id") %>%
  left_join(select(tiendas, tienda_id, nombre_tienda), by = "tienda_id")
stopifnot(nrow(maestra_t) == nrow(transacciones))   # control

resumen <- maestra_t %>%
  group_by(nombre_tienda) %>%
  summarise(ingresos = sum(total_transaccion), transacciones = n(),
            ticket_promedio = round(mean(total_transaccion), 2),
            clientes_unicos = n_distinct(cliente_id), .groups = "drop") %>%
  mutate(ranking = min_rank(desc(ingresos))) %>%
  arrange(ranking)

trimestres <- maestra_t %>%
  group_by(nombre_tienda, trimestre) %>%
  summarise(ingresos = round(sum(total_transaccion)), .groups = "drop") %>%
  pivot_wider(names_from = trimestre, values_from = ingresos,
              values_fill = 0)

top_producto <- maestra_t %>%
  group_by(nombre_tienda, nombre_producto, categoria) %>%
  summarise(ingresos = sum(total_transaccion), .groups = "drop") %>%
  group_by(nombre_tienda) %>%
  slice_max(ingresos, n = 1, with_ties = FALSE) %>%
  ungroup()

dir.create("resultados", showWarnings = FALSE)
write_xlsx(list("Resumen" = resumen, "Trimestres" = trimestres,
                "Top producto" = top_producto),
           "resultados/reporte_tiendas.xlsx")
readxl::excel_sheets("resultados/reporte_tiendas.xlsx")
resumen %>% select(ranking, nombre_tienda, ingresos, ticket_promedio)
top_producto %>% head(3)
\end{rcode}
\begin{rsalida}
[1] "Resumen"      "Trimestres"   "Top producto"
# A tibble: 10 x 4
   ranking nombre_tienda    ingresos ticket_promedio
     <int> <chr>               <dbl>           <dbl>
 1       1 Sucursal Outlet   114377.           1079.
 2       2 Sucursal Sur      105305.            966.
 3       3 Sucursal Plaza A  105216.            965.
 4       4 Sucursal Express  102791.            952.
 5       5 Sucursal Premium  101538.            940.
 6       6 Sucursal Norte     89024.            899.
 7       7 Sucursal Este      88420.            941.
 8       8 Sucursal Centro    88414.           1016.
 9       9 Sucursal Oeste     84672.            891.
10      10 Sucursal Plaza B   80236.            944.
# A tibble: 3 x 4
  nombre_tienda    nombre_producto categoria   ingresos
  <chr>            <chr>           <chr>          <dbl>
1 Sucursal Centro  Producto 40     Electrónica    5611.
2 Sucursal Este    Producto 47     Automotriz     6652.
3 Sucursal Express Producto 27     Juguetes       6072.
\end{rsalida}

El archivo tiene las tres hojas. La Sucursal Outlet vuelve a ser la
número uno (\$114,377 en transacciones), igual que en \codigo{ventas}. Usamos
\codigo{with\_ties = FALSE} en \codigo{slice\_max()} para garantizar un solo
producto por tienda aunque hubiera empates.
\end{solucion}
